\documentclass[11pt,a4paper]{article}
\usepackage{amsmath}
\usepackage{siunitx}
\usepackage[margin=1.8cm]{geometry}
\sisetup{per-mode=symbol}

\begin{document}

\title{Momentum Worksheet}
\author{}
\date{}
\maketitle

\begin{enumerate}
    \item A particle of mass \SI{4.0}{\kilogram} is moving to the right at \SI{3.0}{\meter\per\second}. \\
    Calculate its momentum.

    \item A particle of mass \SI{3.0}{\kilogram} is moving at \SI{5.0}{\meter\per\second} to the right. It collides with a stationary particle of mass \SI{1.0}{\kilogram}. The particles stick together after the collision. \\
    Find the speed of the combined mass immediately after the collision.

    \item A particle of mass \SI{2.0}{\kilogram} is moving to the right with unknown speed $u$. It collides head-on with a particle of mass \SI{3.0}{\kilogram} moving to the left at \SI{4.0}{\meter\per\second}. The particles stick together and move with a speed of \SI{1.0}{\meter\per\second} to the left after the collision. \\
    Find the initial speed $u$ of the \SI{2.0}{\kilogram} particle.

    \item A particle $A$ of mass \SI{0.20}{\kilogram} moves to the right at \SI{6.0}{\meter\per\second} and collides with an identical particle $B$ at rest. After the collision, $A$ moves to the left at \SI{2.0}{\meter\per\second}. \\
    Find the speed and direction of $B$ immediately after the collision.

    \item A particle $P$ of mass \SI{0.50}{\kilogram} moves to the right at \SI{4.0}{\meter\per\second}. It collides head-on with a particle $Q$ of mass \SI{1.0}{\kilogram} moving to the left at \SI{2.0}{\meter\per\second}. After the collision the two particles stick together. \\
    Find the speed and direction of the combined mass after the collision.

    \item A particle of mass \SI{3.0}{\kilogram} is initially at rest. It explodes into two fragments of mass \SI{1.0}{\kilogram} and \SI{2.0}{\kilogram} which move in opposite directions along a straight line. The \SI{1.0}{\kilogram} fragment moves to the right at \SI{6.0}{\meter\per\second}. \\
    Find the velocity of the \SI{2.0}{\kilogram} fragment.

    \item A particle of mass $m$ moves along a straight line with speed $u$. It collides with a stationary particle of mass $2m$. After the collision, the two particles stick together and move as one body. \\
    Find the speed of the combined mass in terms of $u$.

    \item A particle $A$ of mass $m$ moves to the right with speed $u$. It collides with a stationary particle $B$, also of mass $m$. After the collision, $A$ moves to the left with speed $\dfrac{u}{4}$. \\
    Find the speed of $B$ after the collision in terms of $u$.

    \item A particle of mass $3m$ is initially at rest. It explodes into two fragments of mass $m$ and $2m$ that move in opposite directions along a straight line. The fragment of mass $m$ moves to the left with speed $3u$. \\
    Find the speed and direction of the fragment of mass $2m$ in terms of $u$.

    \item A particle $X$ of mass \SI{0.15}{\kilogram} moves to the right at \SI{12}{\meter\per\second}. It collides head-on with a particle $Y$ of mass \SI{0.25}{\kilogram} moving to the left at \SI{8.0}{\meter\per\second}. After the collision, $X$ moves to the left at \SI{4.0}{\meter\per\second}. \\
    Find the speed and direction of $Y$ after the collision.

    \item A particle of mass \SI{0.80}{\kilogram} is at rest. It explodes into two fragments of mass \SI{0.30}{\kilogram} and \SI{0.50}{\kilogram}. The \SI{0.30}{\kilogram} fragment moves to the east at \SI{6.0}{\meter\per\second}. Assuming the motion is along a straight line east--west, \\
    find the speed and direction of the \SI{0.50}{\kilogram} fragment.
\end{enumerate}

\newpage

\section*{Answer Key}

\begin{enumerate}
    \item $p = mv = \SI{4.0}{\kilogram} \times \SI{3.0}{\meter\per\second} = \SI{12}{\kilogram\meter\per\second}$ (to the right).

    \item Total initial momentum $= (3.0)(5.0) = \SI{15}{\kilogram\meter\per\second}$. \\
    Total mass after collision $= 3.0 + 1.0 = \SI{4.0}{\kilogram}$. \\
    \[
        v = \frac{15}{4} = \SI{3.75}{\meter\per\second} \text{ to the right.}
    \]

    \item Momentum before: $2u + 3(-4.0)$. \\
    Momentum after: $(2 + 3)(-1.0) = 5(-1.0) = -5$. \\
    \[
        2u - 12 = -5 \quad\Rightarrow\quad 2u = 7 \quad\Rightarrow\quad u = \SI{3.5}{\meter\per\second} \text{ (to the right).}
    \]

    \item Initial momentum $= (0.20)(6.0) = \SI{1.2}{\kilogram\meter\per\second}$. \\
    After collision:
    \[
        (0.20)(-2.0) + (0.20)v_B = 1.2
    \]
    \[
        -0.40 + 0.20 v_B = 1.2 \quad\Rightarrow\quad 0.20 v_B = 1.6 \quad\Rightarrow\quad v_B = \SI{8.0}{\meter\per\second} \text{ to the right.}
    \]

    \item Initial momentum:
    \[
        (0.50)(4.0) + (1.0)(-2.0) = 2.0 - 2.0 = \SI{0}{\kilogram\meter\per\second}.
    \]
    Total mass after $= 0.50 + 1.0 = \SI{1.5}{\kilogram}$. \\
    \[
        v = \frac{0}{1.5} = \SI{0}{\meter\per\second}.
    \]
    The combined mass is at rest.

    \item Initial momentum $= 0$. After explosion:
    \[
        (1.0)(6.0) + (2.0)v = 0
    \]
    \[
        6.0 + 2.0 v = 0 \quad\Rightarrow\quad v = \SI{-3.0}{\meter\per\second}.
    \]
    The \SI{2.0}{\kilogram} fragment moves at \SI{3.0}{\meter\per\second} to the left.

    \item Initial momentum $= mu + 2m \cdot 0 = mu$. \\
    Total mass after $= m + 2m = 3m$. \\
    \[
        v = \frac{mu}{3m} = \frac{u}{3}.
    \]

    \item Initial momentum $= mu + m\cdot 0 = mu$. After collision:
    \[
        m\left(-\frac{u}{4}\right) + m v_B = mu
    \]
    \[
        -\frac{mu}{4} + m v_B = mu \quad\Rightarrow\quad v_B = mu + \frac{mu}{4}\,\bigg/\,m = \frac{5u}{4}.
    \]
    So $v_B = \dfrac{5u}{4}$ to the right.

    \item Initial momentum $= 0$. After explosion:
    \[
        m(-3u) + 2m v = 0
    \]
    \[
        -3mu + 2m v = 0 \quad\Rightarrow\quad 2m v = 3mu \quad\Rightarrow\quad v = \frac{3u}{2}.
    \]
    The fragment of mass $2m$ moves to the right with speed $\dfrac{3u}{2}$.

    \item Initial momentum:
    \[
        (0.15)(12) + (0.25)(-8.0) = 1.8 - 2.0 = -\SI{0.20}{\kilogram\meter\per\second}.
    \]
    After collision:
    \[
        (0.15)(-4.0) + (0.25)v_Y = -0.20
    \]
    \[
        -0.60 + 0.25 v_Y = -0.20 \quad\Rightarrow\quad 0.25 v_Y = 0.40 \quad\Rightarrow\quad v_Y = \SI{1.6}{\meter\per\second} \text{ to the right.}
    \]

    \item Initial momentum $= 0$. After explosion:
    \[
        (0.30)(6.0) + (0.50) v = 0
    \]
    \[
        1.8 + 0.50 v = 0 \quad\Rightarrow\quad v = \SI{-3.6}{\meter\per\second}.
    \]
    The \SI{0.50}{\kilogram} fragment moves at \SI{3.6}{\meter\per\second} to the west.
\end{enumerate}

\end{document}

